\section{Experimental evaluation}
\label{sec:evaluation}

We evaluate \PurePy empirically as a target language for generated code, with a language model as the code
generator. The evaluation asks four research questions.
\begin{description}
\item[RQ1] Does Python syntax help a model write correct programs?
\item[RQ2] Are the restrictions of \PurePy a barrier to generation?
\item[RQ3a] Do the restrictions of \PurePy improve the correctness of generated programs?
\item[RQ3b] Do the restrictions of \PurePy make generated programs shorter?
\item[RQ4] Can a model translate existing Python into \PurePy?
\end{description}
\secref{evaluation-setup} describes the setup shared by the four questions, and
\secref{evaluation-syntax} to \secref{evaluation-translating} take the questions in turn.
\todo{The research questions are in draft form.}

\subsection{Experimental setup}
\label{sec:evaluation-setup}

A model is given the same problems under several conditions, each with equivalent documentation: \PurePy,
unchecked \PurePy, Python, OCaml, and translation of the benchmark's reference solution into \PurePy.
Unchecked \PurePy keeps the grammar of \PurePy, with annotations optional, and drops the well-formedness
rules, so an answer is checked for syntax only and run by CPython. OCaml is restricted to its pure subset,
with no references or mutable fields, so that it differs from \PurePy mainly in syntax.

For each problem, the model is prompted, the answer is checked statically and run against the benchmark's
official tests, and the number of rounds is recorded. Rejections by the checker are fed back to the model for a
bounded number of repair rounds, and repair stops at acceptance or at the budget. Test outcomes are not fed
back, so a pass measures generation from the problem statement and the checker's diagnostics alone. A \PurePy
answer is a Python program, so the benchmark's Python tests run unchanged against it under CPython, and the
benchmarks are usable as they are.

The benchmarks are HumanEval+ and MBPP+~\cite{liu2023}, which extend HumanEval~\cite{chen2021} and
MBPP~\cite{austin2021} with many more tests per problem, and LiveCodeBench~\cite{jain2025}. HumanEval problems
are small functions over the standard library, so most fall within the \PurePy subset.
MultiPL-E~\cite{cassano2023} provides HumanEval and MBPP translated into OCaml. The self-repair scenario of
LiveCodeBench, in which the model fixes a solution given failing test output, is the nearest precedent for the
repair loop.

The evaluation pipeline in \texttt{evaluation/} implements the \PurePy and Python conditions. It gives a model
the same problems in both, runs the \PurePy checker and the benchmark's official tests on each answer, feeds
checker rejections back to the model for a bounded number of repair rounds, and produces
\figref{evaluation-generation} from the recorded outcomes. The figure reports how many of the selected problems
pass the official tests, and how many are both accepted by the checker and pass the tests, for the initial
answer and after repair. \todo{The committed result file is the pipeline's smoke run on two problems with a
small local model; replace it with the full experiment.} \todo{The unchecked \PurePy, OCaml and translation
conditions are not yet implemented, and nor is the classification of rejections by rule.}

\begin{figure}[t]
\centering
\includegraphics[width=\linewidth]{paper/fig/evaluation-generation}
\caption{Outcomes of the generation experiment, produced by the evaluation pipeline. Left: problems whose final
program passes the official tests. Right: problems whose program is also accepted by the \PurePy checker; the Python
condition is shown by its test outcome only, as it is not checked.}
\label{fig:evaluation-generation}
\end{figure}

\subsection{Python syntax}
\label{sec:evaluation-syntax}

RQ1 asks whether Python syntax helps a model write correct programs. Language models generate Python more
accurately than languages with less training data. MultiPL-T~\cite{cassano2024} measures how much worse models
do in OCaml and Racket, and closes part of the gap by fine-tuning on Python programs translated into those
languages. Because \PurePy programs are Python programs, a model trained on Python is already trained on
\PurePy, and RQ1 asks whether this spares \PurePy the gap without fine-tuning. We compare \PurePy with OCaml
on how often the static check accepts the first attempt, how often the final program passes the tests, and
the number of rounds to an accepted program. A clear gap in \PurePy's favour would support the claim; a small
gap would mean that the argument needs rethinking.

\subsection{Cost of the restrictions}
\label{sec:evaluation-cost}

RQ2 asks whether the restrictions of \PurePy are a barrier to generation. Over all rounds, we count
rejections per forbidden feature. Each diagnostic of the checker is named after the rule it enforces, so a
rejection is classified by its diagnostic. For each rejection, we record whether the next round clears the
feature, and the number of rounds to acceptance. The restrictions are not a barrier if violations are rare and
are cleared in a round or two from the diagnostic alone.

\subsection{Effect of the restrictions on the program}
\label{sec:evaluation-effect}

RQ3a asks whether the restrictions of \PurePy improve the correctness of generated programs. We compare the
pass rate of the final program in \PurePy and in Python; the gap is the effect of the restrictions as a whole.
Comparing \PurePy with unchecked \PurePy measures the effect of the well-formedness rules, and comparing
unchecked \PurePy with Python measures the effect of the grammar. Types may catch errors otherwise revealed
only by the tests, and the model may write better code when it must state types. RQ2 measures the cost of
the restrictions in rejections and rounds, and RQ3a their benefit in correctness.
\figref{evaluation-generation} reports the pass rates of the \PurePy and Python conditions.

RQ3b asks whether the restrictions make generated programs shorter. A program confined to a functional style
may be more concise, although a shorter program is not necessarily a better one. We compare the length of
the final \PurePy program, in non-blank lines and in tokens, with the model's Python answer to the same
problem, which holds the author constant, and with the benchmark's reference solution, which HumanEval+ and
MBPP+ provide but LiveCodeBench lacks.

\subsection{Translating Python to \PurePy}
\label{sec:evaluation-translating}

RQ4 asks whether a model can translate existing Python programs into \PurePy. The model is given the
benchmark's reference solution to each problem and asked for a \PurePy version, under the same loop as for
generation: the checker on the translation, the tests under CPython and the rounds recorded. We report how often
translation succeeds. A breakdown of the failed translations by forbidden feature can follow from the
diagnostics on the final attempt. \todo{Unclear whether this question can be answered without libraries being
part of the specification.}
